\section{09-23 主结果：原几何相位的纯符号解析式}
\label{sec:symbolic-ade}

本节取代 09-22 被否决的 transfer 设计作为当前研究路线：先代入完整
$\Li_2$ 公式并尝试纯符号恒等式，再用同一个四面体的直接几何体积。
不使用数值拟合、数值归约或另一上同调代表替换当前相位。
后文旧代数代表保留为历史附录性质的比较，\emph{不是本节的方法或证据}。
本节的精确实现拒绝浮点输入；历史数值接口仅保留兼容性。

已完成的范围是：一般位置完整解析主程序、$2T$ 的十二类实际有理体积、
$2O$ 的 84 类和 $2I$ 的 563 类完整非退化参数表，以及原 E/F/T
退化链的有限解析接口。必须区分完整解析表达与消去全部特殊函数：
后者对 $2O,2I$ 尚未完成；事实上，要求所有类型都是有理体积已被下文反例排除。
原始任务提示文件当前未在 checkout 中找到，本文按本轮用户明确给出的要求执行。

\subsection{三个群元直接给出六个迹与取向}

固定 $g=x_0I+i\boldsymbol{x}\cdot\boldsymbol{\sigma}$，累计顶点为
\[
(I,g_1,g_1g_2,g_1g_2g_3).
\]
本节边序严格为 $(01,02,03,23,13,12)$：
\begin{align}
\ell_1&=\arccos\tfrac12\Tr g_1,&
\ell_2&=\arccos\tfrac12\Tr(g_1g_2),&
\ell_3&=\arccos\tfrac12\Tr(g_1g_2g_3),\\
\ell_4&=\arccos\tfrac12\Tr g_3,&
\ell_5&=\arccos\tfrac12\Tr(g_2g_3),&
\ell_6&=\arccos\tfrac12\Tr g_2.
\end{align}
对边为 $j$ 与 $j+3$（模 6）。令 $A=g_1,B=g_1g_2,C=g_1g_2g_3$，则
\begin{equation}
 D=\det(q(I),q(A),q(B),q(C))=\tfrac14\Tr\bigl(A(BC-CB)\bigr).
 \label{eq:exact-orientation}
\end{equation}
由 $BC-CB=-2i(\boldsymbol b\times\boldsymbol c)\cdot\boldsymbol\sigma$
直接得右边为 $\boldsymbol a\cdot(\boldsymbol b\times\boldsymbol c)$。
不可把 $A,B,C$ 误写成三个增量本身。输入
$e_1,e_2,(0,3/5,0,4/5)$ 时正确值为 $3/5$，误用增量会得到 $4/5$。

$D\ne0$ 保证顶点线性无关、正锥不含零，从而存在共同开半球，
对应唯一凸短测地线四面体。单四面体入口拒绝 $D=0$，绝不以
$\operatorname{sgn}(0)=0$ 把全局相位置为 1；例如循环子群全部位于大圆，
仍必须保留其进位公式或全局链构造。

\subsection{八个二重对数及六个显式对数项}

以下是 Murakami 定理 1.2 的边长代入及固定 $z$ 偏导展开，
没有隐式积分或待求偏导。记 $c_j=\cos\ell_j$，
$b_j=c_j+i\sqrt{1-c_j^2}$，$a_j=-b_{j+3}^{-1}$，正平方根及主支角约定不变。
\begin{align}
q_0={}&a_1a_4+a_2a_5+a_3a_6+a_1a_2a_6+a_1a_3a_5\notag\\
 &+a_2a_3a_4+a_4a_5a_6+\prod_{j=1}^6a_j,\\
q_1={}&-\sum_{j=1}^3(a_j-a_j^{-1})(a_{j+3}-a_{j+3}^{-1}),
\qquad q_2=\overline{q_0},\\
\Delta={}&q_1^2-4q_0q_2=16\det G,\qquad G=V^{\mathsf T}V,\\
z={}&\frac{-q_1+\sqrt\Delta}{2q_2}
     =\frac{-2q_0}{q_1+4\sqrt{\det G}}.
\end{align}
其中 $q_2$ 也等于 $q_0$ 每个单项式取倒数后的和。为完整指定八个变元，令
\begin{align}
 Z_0&=z,&Z_1&=b_1b_2b_4b_5z,\\
 Z_2&=b_1b_3b_4b_6z,&Z_3&=b_2b_3b_5b_6z,\\
 Z_4&=b_4b_5b_6z,&Z_5&=b_2b_3b_4z,\\
 Z_6&=b_1b_3b_5z,&Z_7&=b_1b_2b_6z.
\end{align}
记 $L_r=\log(1-Z_r)$（各自取主支，不能擅自合并）：
\begin{align}
\Sigma_1&=L_1+L_2-L_6-L_7,&\Sigma_2&=L_1+L_3-L_5-L_7,\\
\Sigma_3&=L_2+L_3-L_5-L_6,&\Sigma_4&=L_1+L_2-L_4-L_5,\\
\Sigma_5&=L_1+L_3-L_4-L_6,&\Sigma_6&=L_2+L_3-L_4-L_7.
\end{align}
完整主程序计算
\begin{align}
W={}&\tfrac12\Re\left[\sum_{r=0}^3\Li_2(Z_r)-\sum_{r=4}^7\Li_2(Z_r)\right]
 -\tfrac12\sum_{j=1}^3(\pi-\ell_j)(\pi-\ell_{j+3})\notag\\
&-\pi\Arg(-q_2)
 -\tfrac12\sum_{j=1}^6\ell_j\bigl(\pi-\ell_{j+3}+\Im\Sigma_j\bigr)-\tfrac{\pi^2}{2}.
\end{align}
\begin{equation}
\boxed{\omega_k(g_1,g_2,g_3)={}
 \exp\!\left[-\frac{ik}{\pi}\operatorname{sgn}(D)W\right]},\qquad k\in\mathbb Z.
\label{eq:exact-master-phase}
\end{equation}
$W$ 是体积模 $2\pi^2$；实际正体积为其在 $(0,\pi^2)$ 的代表。
整数 level 的相位不需要通过小数来选择代表。

\paragraph{主支控制。}
$q_1=4\sum\sin\ell_j\sin\ell_{j+3}>0$，$\Delta>0$，故稳定根满足
$|z|<1$，全部 $Z_r$ 同模；因而 $1-Z_r$ 的实部为正，避免了对数支割。
正定相关矩阵域连通，Murakami 的驻点恒等式
$\exp(2z\partial_zL_e)=1$ 使连续量 $z\partial_zL_e$ 为常数 $\pi i$
（在正交例校准）；若 $q_0=0$，稳定根为零，导数乘积将为零，矛盾。
$\Arg(-q_2)$ 的跳跃只改变整球体积。
这里需要实际 Gram 正定，不能用任意六个数满足 $\det G>0$ 来代替。

\subsection{已经完成的严格符号化简}

通用约简器仅在已证明的域内使用共轭、Euler 反射及倍角恒等式：
\begin{align}
\Re\Li_2(\bar z)&=\Re\Li_2(z),&& |z|<1,\\
\Li_2(z)+\Li_2(1-z)&=\pi^2/6-\log z\log(1-z),&&|z|,|1-z|<1,\\
\Li_2(z)+\Li_2(-z)&=\tfrac12\Li_2(z^2),&&|z|<1.
\end{align}
未知域不强行约简，剩余项明确返回；这不是所有多重对数恒等式的完备判定器。

\paragraph{正交校准。}
输入 $(e_1,e_3,e_1)$ 给累计顶点 $(e_0,e_1,e_2,e_3)$，
$q_0=-4-4i,q_1=12,q_2=-4+4i,z=(1+i)/2$。
四个正项与四个负项互为共轭，故 $\Li_2$ 实部完全抵消。
又 $\Sigma_j=-i\pi$、$\Arg(-q_2)=-\pi/4$，所以
\[
 W=-\frac{3\pi^2}{8}+\frac{\pi^2}{4}+\frac{3\pi^2}{4}-\frac{\pi^2}{2}
   =\frac{\pi^2}{8},\qquad \omega_k=e^{-ik\pi/8}.
\]
这也等于 $S^3$ 正交卦限的体积 $2\pi^2/16$。

\paragraph{两参数族的完整 $\Li_2$ 恒等式证明。}
令边长为 $(\alpha,\pi/2,\pi/2,\beta,\pi/2,\pi/2)$，
$0<\alpha,\beta<\pi$。设 $p=e^{i\alpha},q=e^{i\beta}$，
$S=pq+p+q-1$，$\gamma=\Arg S$，$\eta=\gamma-(\alpha+\beta)/2$。
由于
\[
 S=2e^{i(\alpha+\beta)/2}
 \left(\cos\frac{\alpha-\beta}{2}+i\sin\frac{\alpha+\beta}{2}\right),
\]
有 $0<\eta<\pi/2$、$0<\gamma<\pi$，且 $z=-2/S$。八项合为两倍
\[
F=\Li_2(z)+\Li_2(-pqz)-\Li_2(-pz)-\Li_2(-qz).
\]
分解
\begin{align}
1-z&=(p+1)(q+1)/S,&1+pqz&=-(p-1)(q-1)/S,\\
1+pz&=(p+1)(q-1)/S,&1+qz&=(p-1)(q+1)/S.
\end{align}
记这四个主支对数分别为 $L_U,L_V,L_X,L_Y$。
逐项计算模和主辐角（正交点的符号为 $-i\pi$）给出
\begin{align}
 L_U+L_V-L_X-L_Y&=-i\pi,\\
 L_X-L_V&=\log\cot(\alpha/2)+i\pi/2,\\
 L_Y-L_V&=\log\cot(\beta/2)+i\pi/2.
\end{align}
用 $d\Li_2(w)=-\log(1-w)d\log w$ 得
\[
dF=i\pi\,d\log z+i(L_X-L_V)d\alpha+i(L_Y-L_V)d\beta,
\quad d\Re F=\pi\,d\eta.
\]
正交点 $\eta=\pi/4,\Re F=0$ 固定积分常数，故
$\Re F=\pi\eta-\pi^2/4$。
六个 $\Im\Sigma_j=-\pi$，$\Arg(-q_2)=\gamma-\pi$，代回主式即
\[
 W=\alpha\beta/2.
\]
允许实数 $-\pi<\alpha<0$ 时，边长用 $|\alpha|$，取向使
$V_{\rm or}=\alpha\beta/2$。独立参数化
\[
x=(\cos\rho\cos u,\cos\rho\sin u,\sin\rho\cos v,\sin\rho\sin v)
\]
的有向体积形式为 $\cos\rho\sin\rho\,du\wedge d\rho\wedge dv$，
$u:0\to\alpha,\rho:0\to\pi/2,v:0\to\beta$，积分也为 $\alpha\beta/2$。
特别地 $(\pm\pi/3,\pi/4)$ 给出 $\pm\pi^2/24$；相同边长仍需取向信息。

\section{2T 的十二类实际体积及几何证明}
\label{sec:2t-actual}

$2T$ 点集为 $\{\pm e_j\}$ 与全部 $(\pm1,\pm1,\pm1,\pm1)/2$。
枚举 $\binom{24}{4}=10626$ 个无序四点集合，5586 个秩不足 4，5040 个非退化。
六个 Gram 元在 24 种顶点排列下取规范键，恰有十二类；满秩矩阵具有相同 Gram
当且仅当由一个 $O(4)$ 等距变换联系，因此此分类确实是无向几何分类。

表中 $K=2(G_{01},G_{02},G_{03},G_{12},G_{13},G_{23})$ 的置换规范键，
$m$ 为下文严格证明的室数。
\begin{center}\small
\begin{tabular}{@{}rrrrrr|rrc@{}}
\toprule
\multicolumn{6}{c|}{$K$}&重数&$m$&$\Vol/\pi^2$\\\midrule
-1&-1&-1&0&0&0&192&270&$15/32$\\
-1&-1&-1&0&1&1&576&84&$7/48$\\
-1&-1&0&0&-1&1&576&120&$5/24$\\
-1&-1&0&1&-1&-1&288&204&$17/48$\\
-1&-1&0&1&1&1&576&36&$1/16$\\
-1&-1&1&0&0&0&576&66&$11/96$\\
-1&-1&1&1&-1&0&576&60&$5/48$\\
-1&0&0&1&1&0&576&30&$5/96$\\
-1&0&1&1&0&1&576&24&$1/24$\\
0&0&0&0&0&0&48&72&$1/8$\\
0&0&1&0&1&1&192&18&$1/32$\\
0&1&1&1&1&1&288&12&$1/48$\\\bottomrule
\end{tabular}\end{center}

\subsection{为什么有限室计数恰好是实际体积}

$F_4$ 短根为 $\{\pm e_j\}$ 与半符号根，长根为 $\{\pm e_i\pm e_j\}$。
其中心超平面共 24 个：4 个坐标面、12 个双坐标面、8 个全符号面。
带符号置换与
\[
H=\frac12\begin{pmatrix}1&1&1&1\\1&1&-1&-1\\1&-1&1&-1\\1&-1&-1&1\end{pmatrix}
\]
生成保持根集的正交群 $W$；$H$ 正是根 $(1,-1,-1,-1)/2$ 的反射。
$W$ 在 24 个短根上传递，$e_0$ 稳定子必置换其余六个正交短根，恰为阶 48 的
$B_3$。所以 $|W|=24\cdot48=1152$。

任意三个独立 Hurwitz 顶点张成的面都在这些超平面内。证明是先把第一个顶点送到
$e_0$，另两点为轴或半符号点：轴/轴给轴法向；轴/半符号点给双坐标法向；
两个半符号点在剩余三维的立方体投影，独立叉积只有两个等绝对值的非零分量。
变回后仍是 $F_4$ 根法向。

基本室可取
\[
x_1>x_2>x_3>0,\qquad x_0>x_1+x_2+x_3.
\]
它的内部没有任何上述超平面。四个 inward 法向为
$e_1-e_2,e_2-e_3,e_3,(e_0-e_1-e_2-e_3)/2$。
带根长颜色的角图为 $3$--$4$--$3$、颜色依次为长/长/短/短，
保持颜色与邻接的自同构只有恒等，故室稳定子平凡。
反射跨越相邻墙使 $W$ 在所有室上传递，所以有 1152 个全等室，
每室体积为 $\pi^2/576$。四面体的四个面都在墙上，因而它是完整小室的并，
边界体积为零。这不是采样估计。

\subsection{纯整数的群元体积公式}

令 $P$ 为 $(8,3,2,1),(7,4,3,2),(6,5,4,1)$ 的全部带符号排列。
三个不交集合各有 384 个点；第一个点在基本室内，而且
\[
H(8,3,2,1)=(7,4,3,2),\qquad H(8,3,2,-1)=(6,5,4,1).
\]
因此这 1152 点属于同一 $W$ 轨道，恰好每个室一个。对顶点矩阵 $V$，
\begin{align}
m(V)&=\#\{p\in P:(V^{-1}p)_j>0\text{ 对所有 }j\},\\
V_{\rm or}&=\operatorname{sgn}(\det V)\frac{m(V)\pi^2}{576},\\
\omega_k&=\exp\left[-\frac{ik\pi}{576}\operatorname{sgn}(\det V)m(V)\right].
\end{align}
全部正负号由整数余子式判定。这是原四面体的实际体积，不包含 $\beta$ 或代表转换。
各类型先尝试 $\Li_2$ 代入，未能消去的项由此几何恒等式识别。

\section{2O 与 2I：完整解析分类与有理体积的障碍}
\label{sec:oi-exact}

$2O=2T\cup\{(\pm e_i\pm e_j)/\sqrt2\}$。
$2I=2T$ 加上 $(0,\pm1/2,\pm\varphi/2,\pm\varphi^{-1}/2)$ 的所有偶坐标置换，
其中 $\varphi=(1+\sqrt5)/2$。
坐标分别属于 $\mathbb Q(\sqrt2)$、$\mathbb Q(\sqrt5)$，程序使用整数对，
不靠浮点阈值判定秩或比较 Gram 键。

左乘等距变换把一个顶点移到 $e_0$，枚举剩余三个不同点，再在 24 个顶点排列下
规范化 Gram。每个无序集合有四种被选中的基点，故全部集合重数为锚定重数乘
$|G|/4$，有序三输入重数为锚定重数乘 6。
\begin{center}\small
\begin{tabular}{@{}crrrrr@{}}\toprule
群&类型&锚定非退化&锚定退化&全部非退化集合&非退化三输入\\\midrule
$2T$&12&840&931&5040&5040\\
$2O$&84&10080&6135&120960&60480\\
$2I$&563&212520&61299&6375600&1275120\\\bottomrule
\end{tabular}\end{center}

完整电子参数表为 \path{results/SU2_symbolic_ADE_exact.json}：每行含代表坐标索引、
全部六个 Gram 元、重数、公式次序的六个余弦及 $\det G$。
共同的 \texttt{formula\_template} 是显式顺序赋值程序，每行
\texttt{formula\_inputs} 给齐代入参数，直至式 \eqref{eq:exact-master-phase}，
不依赖另一个体积数据库或未展示的相位公式。
锚定代表与规范 Gram 可能有不同顶点次序，但 $\det G$ 和无向体积不变；
有序输入自己的 $\operatorname{sgn}D$ 必须另行恢复。
三群共 659 类型逐类记录符号化简尝试；时间预算耗尽仅表示该次 CAS 尝试未结束，
绝不是不可化简证明。

\subsection{2O 中实际无理体积的严格反例}

取累计顶点
\[
e_0,\quad e_1,\quad (e_1+e_2)/\sqrt2,\quad(e_2+e_3)/\sqrt2,
\qquad D=1/2.
\]
对应三个输入是 $e_1,(e_0+e_3)/\sqrt2,(e_0+e_1-e_2+e_3)/2$，均属于 $2O$。
后三点所在 $S^2$ 三角形的内角为
$\pi/4,\pi-\arctan\sqrt2,\pi/2-\arctan\sqrt2$。
Girard 公式给面积 $\Omega=3\pi/4-2\arctan\sqrt2$。
设 $\theta=\Omega/2$，则
$\tan\theta=1/(3+\sqrt2)$、$0<\theta<\pi/8$。
与正交点 $e_0$ 作球面 join，积分 $\int_0^{\pi/2}\sin^2u\,du=\pi/4$，所以
\begin{equation}
\boxed{V_{2O}=\frac\pi2\arctan\frac1{3+\sqrt2}},\qquad
\omega_k=\exp\left[-\frac{ik}{2}\arctan\frac1{3+\sqrt2}\right].
\end{equation}
令 $t=\tan\theta$，有
$e^{2i\theta}=(1+it)/(1-it)\in\mathbb Q(\sqrt2,i)=\mathbb Q(\zeta_8)$。
这个域的单位根只有 $\mu_8$，但 $0<2\theta<\pi/4$ 排除所有这些根。
所以 $\theta/\pi$、$V_{2O}/\pi^2$ 均无理，非零整数 level 的相位也不是单位根。

\subsection{2I 中实际无理体积的严格反例}

取累计顶点
\[
 e_0,e_1,e_2,(0,1/2,\varphi/2,\varphi^{-1}/2),\qquad D=\varphi^{-1}/2>0.
\]
对应输入 $e_1,e_3,(\varphi/2,\varphi^{-1}/2,0,-1/2)$ 均属于 $2I$。
空间三角形的 Girard 公式及半角恒等式给出
\[
\tan\frac\Omega2=\frac{\det(A,B,C)}{1+A\cdot B+B\cdot C+C\cdot A}
 =\frac1{4\varphi+1}.
\]
同一个 join 积分产生
\begin{equation}
\boxed{V_{2I}=\frac\pi2\arctan\frac1{4\varphi+1}},\qquad
\omega_k=\exp\left[-\frac{ik}{2}\arctan\frac1{4\varphi+1}\right].
\end{equation}
此时 $e^{2i\theta}\in\mathbb Q(\sqrt5,i)$，其单位根只有 $\mu_4$。
具体地，圆分次数必须整除 4；$\mu_3,\mu_6$ 要求 $\sqrt{-3}$，
$\mu_8$ 要求 $\sqrt2$，$\mu_{12}$ 要求 $\sqrt3$，都不在此域内。
$\mu_5,\mu_{10}$ 给循环四次域，不能等于双二次域。
而 $0<2\theta<\pi/2$ 排除 $\mu_4$，所以同样得到无理体积及非单位根相位。
$2O$ 的 $\mu_8$ 判断也可由同一低次数圆分域列表核验。

这两个反例说明：\emph{通过坐标符号把 $2O,2I$ 所有类型都化成有理体积并不可行}。
它不妨碍完整 $\Li_2$/对数解析式，也不妨碍个别类型具有更简短的几何表达；
不得把符号约简尚未完成与已证明的无理性混为一谈。
这不与有限群反常类的有限阶矛盾：类的 $N$ 次幂平凡只表示
$\omega^N$ 是某个余边界，不意味着原代表每个逐点值都满足 $\omega^N=1$。
本轮并未利用这种转换更换代表。

\section{精确退化覆盖、复现与本轮边界}

全局接口沿用原 E/F/T 代表：相邻重复规范化，精确判定零是否在顶点凸包，
不安全边经固定 $J=e_1$ 分段，不安全面及三链按原规则锥接边界。
锥点取首个不落在相关张成空间内的 $a(1,t,t^2,t^3)$，$0\le t\le3M$。
每个真子空间至多排除三个整数候选，所以过程有限，每面至多 6 项，每三链至多 24 项。
安全的降秩项积分为零但保留链记录；满秩项按正范数单位化后代入本节完整解析式。
最后取各项系数乘定向体积的和，再求负号相位。
链边界相容和闭三链整球体积证明五边形，不能仅由有限测试声称所有输入成立。

锥点通常不属于 E 群，不能强行归入十二、84 或 563 类型表。
$2T$ 非退化群顶点的有理体积表，不等于证明了全部退化填充单形也有有理体积。
\texttt{expand=False} 仅返回链审计；有未展开满秩项时相位字段为
\texttt{None}。默认展开给精确特殊函数表达式，\texttt{reduce=True} 才请求额外 CAS 化简。

\begin{itemize}
\item \path{src/su2_symbolic.py}：群元/矩阵入口、完整主式、保守符号约简及双圆弧证书。
\item \path{src/su2_2t_geometry.py}：十二类及严格 $F_4$ 几何室公式。
\item \path{src/su2_exact_ade.py}：三群精确分类、无理体积证书及原退化链解析式。
\item \path{scripts/check_symbolic_2t.py}：逐类型符号尝试、公共公式模板及全参数导出。
\item \path{docs/review-0923/SU2_LI2_ADE_EXACT.md}：完整文字推导、命令与审核记录。
\end{itemize}
复现新数学模块只需 SymPy；不能把后文历史浮点误差当成本节的证明。
完整回归为 70 项通过，其中 31 项为新增纯精确测试。
659 个类型均有批量记录；每类四秒的符号尝试结果如下：
\begin{center}
\begin{tabular}{@{}crrrr@{}}\toprule
群&消去 $\Li_2$&明确剩余项&预算内未完成&总数\\\midrule
$2T$&3&9&0&12\\
$2O$&17&42&25&84\\
$2I$&33&40&490&563\\\bottomrule
\end{tabular}\end{center}
消去 $\Li_2$ 不保证全部 $\log/\Arg$ 也被消去，更不等同于有理体积；
预算内未完成不等同于不可化简。完整解析模板与参数覆盖不受预算影响。
本报告可供专家提出意见，不声称已经通过外部同行评议。
% Generated by python -m scripts.render_exact_catalogue; exact integer enumeration.
\section{2O、2I 全部非退化类型：参数与精确体积短表}
\label{sec:complete-gram-table}
本节在 PDF 内完整列出所有类型及其精确体积，不要求专家另行运行代码。
每个六位键 $K=d_1d_2d_3d_4d_5d_6$ 后的下标是锚定重数 $m$。
从各群的字母表取 $x_j=A_{d_j}$（字母索引从 0 开始），构造
\[
G(K)=\begin{pmatrix}1&x_1&x_2&x_3\\x_1&1&x_4&x_5\\
x_2&x_4&1&x_6\\x_3&x_5&x_6&1\end{pmatrix},\qquad
(c_1,\ldots,c_6)=(x_1,x_2,x_3,x_6,x_5,x_4).
\]
这些 $c_j$ 与 $\det G(K)$ 可代入第 \ref{sec:symbolic-ade} 节的完整相位式；
同一类型的实际体积由第 \ref{sec:polar-dual} 节的极对偶公式给出：
\[
V=a\pi^2+\frac{\pi}{2}\sum_jb_j\Theta_j,\qquad
\omega_k=\exp\!\Bigl[-ik\operatorname{sgn}(D)\Bigl(a\pi+\tfrac12\sum_jb_j\Theta_j\Bigr)\Bigr],
\]
其中两群的基底常数为
\[
2O:\ \Theta=\beta=\arctan\tfrac1{\sqrt2},\qquad
2I:\ \Theta=(A,B,C)=\bigl(\arctan\tfrac12,\arctan\tfrac1{\sqrt5},\arctan\sqrt{15}\bigr).
\]
$N$ 是极对偶中的 Coxeter 室数（$|W|=1152$ 或 $14400$），$a,b_j$ 为有理数；
$b=0$ 当且仅当体积是 $\pi^2$ 的有理倍数（定理 \ref{thm:angle-independence}）。
有序群输入另给 $\operatorname{sgn}D$；只看无向类型不能恢复取向。
字母表按整数对编码排序，并非按实数大小排序；行序与 JSON 的零起始类型索引一致。
\subsection{2O：84 类（其中 60 类为有理体积）}
\[ A=(-1,-\tfrac12,-\tfrac{\sqrt2}{2},0,\tfrac{\sqrt2}{2},\tfrac12,1).\]
{\footnotesize\begin{longtable}{@{}lrll@{\hspace{2em}}lrll@{}}
\toprule
$K_m$ & $N$ & $a$ & $b$ & $K_m$ & $N$ & $a$ & $b$\\\midrule\endhead
\texttt{111333}$_{32}$ & 6 & $\tfrac{15}{32}$ & 0 & \texttt{111355}$_{96}$ & 36 & $\tfrac{7}{48}$ & 0\\
\texttt{112334}$_{192}$ & 15 & $\tfrac{15}{64}$ & 0 & \texttt{112533}$_{96}$ & 6 & $\tfrac{29}{96}$ & $1$\\
\texttt{112544}$_{96}$ & 66 & $\tfrac{7}{96}$ & 0 & \texttt{113315}$_{96}$ & 24 & $\tfrac{5}{24}$ & 0\\
\texttt{113323}$_{192}$ & 3 & $\tfrac{35}{64}$ & 0 & \texttt{113333}$_{96}$ & 24 & $\tfrac{1}{4}$ & 0\\
\texttt{113334}$_{192}$ & 45 & $\tfrac{5}{64}$ & 0 & \texttt{113511}$_{48}$ & 12 & $\tfrac{17}{48}$ & 0\\
\texttt{113523}$_{192}$ & 18 & $\tfrac{19}{96}$ & 0 & \texttt{113533}$_{96}$ & 48 & $\tfrac{1}{4}$ & $-1$\\
\texttt{113534}$_{192}$ & 78 & $\tfrac{5}{96}$ & 0 & \texttt{113555}$_{96}$ & 84 & $\tfrac{1}{16}$ & 0\\
\texttt{114323}$_{192}$ & 33 & $\tfrac{9}{64}$ & 0 & \texttt{114522}$_{96}$ & 30 & $\tfrac{17}{96}$ & 0\\
\texttt{114533}$_{96}$ & 90 & $-\tfrac{5}{96}$ & $1$ & \texttt{115333}$_{96}$ & 42 & $\tfrac{11}{96}$ & 0\\
\texttt{115513}$_{96}$ & 60 & $\tfrac{5}{48}$ & 0 & \texttt{122335}$_{96}$ & 2 & $\tfrac{193}{288}$ & $-1$\\
\texttt{122345}$_{192}$ & 19 & $\tfrac{91}{576}$ & 0 & \texttt{122445}$_{96}$ & 44 & $\tfrac{11}{144}$ & 0\\
\texttt{123325}$_{96}$ & 10 & $\tfrac{77}{288}$ & 0 & \texttt{123331}$_{96}$ & 1 & $\tfrac{385}{576}$ & 0\\
\texttt{123333}$_{192}$ & 12 & $\tfrac{5}{16}$ & 0 & \texttt{123335}$_{192}$ & 23 & $-\tfrac{1}{576}$ & $1$\\
\texttt{123341}$_{96}$ & 14 & $\tfrac{55}{288}$ & 0 & \texttt{123353}$_{192}$ & 21 & $\tfrac{47}{192}$ & $-1$\\
\texttt{123421}$_{96}$ & 5 & $\tfrac{245}{576}$ & 0 & \texttt{123431}$_{192}$ & 22 & $\tfrac{35}{288}$ & 0\\
\texttt{123433}$_{192}$ & 48 & $\tfrac{1}{8}$ & 0 & \texttt{123435}$_{192}$ & 74 & $\tfrac{13}{288}$ & 0\\
\texttt{123445}$_{192}$ & 91 & $\tfrac{19}{576}$ & 0 & \texttt{123453}$_{192}$ & 63 & $\tfrac{13}{192}$ & 0\\
\texttt{124421}$_{48}$ & 52 & $\tfrac{13}{144}$ & 0 & \texttt{124431}$_{96}$ & 77 & $\tfrac{29}{576}$ & 0\\
\texttt{125332}$_{192}$ & 9 & $\tfrac{67}{192}$ & 0 & \texttt{125433}$_{192}$ & 81 & $\tfrac{11}{192}$ & 0\\
\texttt{133331}$_{48}$ & 32 & $\tfrac{2}{9}$ & 0 & \texttt{133335}$_{96}$ & 64 & $\tfrac{1}{9}$ & 0\\
\texttt{133341}$_{96}$ & 73 & $\tfrac{25}{576}$ & 0 & \texttt{133343}$_{192}$ & 84 & $\tfrac{1}{16}$ & 0\\
\texttt{133345}$_{192}$ & 95 & $\tfrac{71}{576}$ & $-1$ & \texttt{133352}$_{192}$ & 51 & $\tfrac{17}{192}$ & 0\\
\texttt{133353}$_{192}$ & 72 & $0$ & $1$ & \texttt{133354}$_{192}$ & 93 & $\tfrac{7}{192}$ & 0\\
\texttt{133445}$_{96}$ & 146 & $-\tfrac{23}{288}$ & $1$ & \texttt{133454}$_{192}$ & 135 & $\tfrac{5}{192}$ & 0\\
\texttt{133553}$_{96}$ & 102 & $\tfrac{5}{96}$ & 0 & \texttt{134435}$_{96}$ & 154 & $\tfrac{5}{288}$ & 0\\
\texttt{134533}$_{192}$ & 123 & $\tfrac{25}{192}$ & $-1$ & \texttt{135535}$_{96}$ & 120 & $\tfrac{1}{24}$ & 0\\
\texttt{222555}$_{32}$ & 18 & $\tfrac{23}{96}$ & 0 & \texttt{223533}$_{96}$ & 24 & $\tfrac{1}{8}$ & $1$\\
\texttt{223535}$_{192}$ & 39 & $\tfrac{7}{64}$ & 0 & \texttt{225522}$_{24}$ & 4 & $\tfrac{73}{144}$ & 0\\
\texttt{225523}$_{96}$ & 29 & $\tfrac{77}{576}$ & 0 & \texttt{225533}$_{96}$ & 46 & $\tfrac{47}{288}$ & $-1$\\
\texttt{233332}$_{48}$ & 18 & $\tfrac{9}{32}$ & 0 & \texttt{233333}$_{96}$ & 36 & $\tfrac{3}{16}$ & 0\\
\texttt{233334}$_{96}$ & 54 & $\tfrac{3}{32}$ & 0 & \texttt{233353}$_{192}$ & 60 & $\tfrac{3}{16}$ & $-1$\\
\texttt{233355}$_{192}$ & 69 & $-\tfrac{3}{64}$ & $1$ & \texttt{233555}$_{96}$ & 102 & $\tfrac{13}{96}$ & $-1$\\
\texttt{235533}$_{96}$ & 97 & $-\tfrac{95}{576}$ & $2$ & \texttt{235534}$_{96}$ & 110 & $\tfrac{7}{288}$ & 0\\
\texttt{333333}$_{32}$ & 72 & $\tfrac{1}{8}$ & 0 & \texttt{333334}$_{96}$ & 108 & $\tfrac{1}{16}$ & 0\\
\texttt{333345}$_{192}$ & 132 & $-\tfrac{1}{16}$ & $1$ & \texttt{333355}$_{96}$ & 120 & $\tfrac{1}{4}$ & $-2$\\
\texttt{333445}$_{96}$ & 168 & $\tfrac{1}{8}$ & $-1$ & \texttt{333555}$_{32}$ & 144 & $-\tfrac{1}{4}$ & $3$\\
\texttt{334433}$_{48}$ & 162 & $\tfrac{1}{32}$ & 0 & \texttt{334535}$_{192}$ & 171 & $-\tfrac{5}{64}$ & $1$\\
\texttt{334555}$_{96}$ & 186 & $\tfrac{11}{96}$ & $-1$ & \texttt{335355}$_{32}$ & 138 & $\tfrac{1}{32}$ & 0\\
\texttt{335445}$_{192}$ & 201 & $\tfrac{1}{64}$ & 0 & \texttt{335533}$_{48}$ & 128 & $\tfrac{1}{18}$ & 0\\
\texttt{335534}$_{96}$ & 169 & $\tfrac{121}{576}$ & $-2$ & \texttt{335544}$_{96}$ & 190 & $-\tfrac{25}{288}$ & $1$\\
\texttt{345544}$_{96}$ & 245 & $\tfrac{5}{576}$ & 0 & \texttt{355555}$_{48}$ & 204 & $\tfrac{1}{48}$ & 0\\
\texttt{444555}$_{32}$ & 270 & $\tfrac{1}{96}$ & 0 & \texttt{445544}$_{24}$ & 292 & $\tfrac{1}{144}$ & 0\\
\bottomrule
\end{longtable}}
\subsection{2I：563 类（其中 268 类为有理体积）}
\[ A=(-1,-\tfrac12,-\tfrac{1+\sqrt5}{4},\tfrac{-1+\sqrt5}{4},0,\tfrac{1-\sqrt5}{4},\tfrac{1+\sqrt5}{4},\tfrac12,1).\]
{\footnotesize\begin{longtable}{@{}lrll@{\hspace{2em}}lrll@{}}
\toprule
$K_m$ & $N$ & $a$ & $b$ & $K_m$ & $N$ & $a$ & $b$\\\midrule\endhead
\texttt{111444}$_{160}$ & 75 & $\tfrac{15}{32}$ & 0 & \texttt{111445}$_{240}$ & 2 & $\tfrac{3091}{3600}$ & $(-\tfrac{1}{2},0,0)$\\
\texttt{111446}$_{240}$ & 382 & $\tfrac{401}{3600}$ & $(\tfrac{1}{2},0,0)$ & \texttt{111447}$_{240}$ & 240 & $\tfrac{3}{10}$ & $(0,1,-\tfrac{1}{2})$\\
\texttt{111456}$_{480}$ & 205 & $\tfrac{221}{1440}$ & 0 & \texttt{111457}$_{480}$ & 103 & $\tfrac{1783}{7200}$ & $(0,0,\tfrac{1}{2})$\\
\texttt{111467}$_{480}$ & 593 & $-\tfrac{487}{7200}$ & $(0,0,\tfrac{1}{2})$ & \texttt{111477}$_{480}$ & 450 & $\tfrac{7}{48}$ & 0\\
\texttt{111557}$_{240}$ & 18 & $\tfrac{853}{1200}$ & $(0,0,-\tfrac{1}{2})$ & \texttt{111577}$_{240}$ & 236 & $\tfrac{509}{1800}$ & $(0,0,-1)$\\
\texttt{111667}$_{240}$ & 1122 & $-\tfrac{83}{1200}$ & $(0,0,\tfrac{1}{2})$ & \texttt{111677}$_{240}$ & 916 & $\tfrac{499}{1800}$ & $(0,0,-1)$\\
\texttt{111777}$_{80}$ & 720 & $-\tfrac{1}{10}$ & $(0,0,1)$ & \texttt{112434}$_{480}$ & 5 & $\tfrac{209}{288}$ & 0\\
\texttt{112436}$_{480}$ & 353 & $\tfrac{893}{7200}$ & 0 & \texttt{112437}$_{480}$ & 181 & $\tfrac{2041}{7200}$ & 0\\
\texttt{112446}$_{480}$ & 117 & $\tfrac{319}{2400}$ & $(-\tfrac{1}{2},0,0)$ & \texttt{112447}$_{480}$ & 39 & $\tfrac{1253}{2400}$ & $(0,-1,0)$\\
\texttt{112467}$_{480}$ & 501 & $\tfrac{49}{800}$ & 0 & \texttt{112533}$_{240}$ & 28 & $\tfrac{997}{1800}$ & 0\\
\texttt{112546}$_{480}$ & 66 & $\tfrac{117}{400}$ & 0 & \texttt{112577}$_{240}$ & 160 & $\tfrac{2}{9}$ & 0\\
\texttt{112637}$_{480}$ & 424 & $\tfrac{83}{900}$ & 0 & \texttt{112644}$_{240}$ & 40 & $\tfrac{143}{360}$ & $(-\tfrac{1}{2},0,0)$\\
\texttt{112666}$_{240}$ & 1052 & $\tfrac{53}{1800}$ & 0 & \texttt{112733}$_{240}$ & 230 & $\tfrac{89}{720}$ & $(0,0,\tfrac{1}{2})$\\
\texttt{112734}$_{480}$ & 72 & $\tfrac{251}{600}$ & $(0,0,-\tfrac{1}{2})$ & \texttt{112736}$_{480}$ & 470 & $\tfrac{113}{720}$ & $(0,0,-\tfrac{1}{2})$\\
\texttt{112744}$_{240}$ & 2 & $\tfrac{2161}{3600}$ & $(0,1,\tfrac{1}{2})$ & \texttt{112747}$_{480}$ & 113 & $\tfrac{233}{7200}$ & $(0,0,\tfrac{1}{2})$\\
\texttt{112766}$_{240}$ & 966 & $\tfrac{151}{1200}$ & $(0,0,-\tfrac{1}{2})$ & \texttt{112767}$_{480}$ & 642 & $-\tfrac{43}{1200}$ & $(0,0,\tfrac{1}{2})$\\
\texttt{112777}$_{240}$ & 436 & $\tfrac{259}{1800}$ & 0 & \texttt{113424}$_{480}$ & 115 & $\tfrac{539}{1440}$ & 0\\
\texttt{113425}$_{480}$ & 17 & $\tfrac{4637}{7200}$ & 0 & \texttt{113427}$_{480}$ & 371 & $\tfrac{551}{7200}$ & 0\\
\texttt{113445}$_{480}$ & 333 & $\tfrac{631}{2400}$ & $(-\tfrac{1}{2},0,0)$ & \texttt{113447}$_{480}$ & 591 & $-\tfrac{43}{2400}$ & $(0,1,0)$\\
\texttt{113457}$_{480}$ & 531 & $\tfrac{79}{800}$ & 0 & \texttt{113527}$_{480}$ & 136 & $\tfrac{227}{900}$ & 0\\
\texttt{113544}$_{240}$ & 200 & $\tfrac{11}{72}$ & $(\tfrac{1}{2},0,0)$ & \texttt{113555}$_{240}$ & 52 & $\tfrac{883}{1800}$ & 0\\
\texttt{113622}$_{240}$ & 172 & $\tfrac{493}{1800}$ & 0 & \texttt{113645}$_{480}$ & 966 & $\tfrac{27}{400}$ & 0\\
\texttt{113677}$_{240}$ & 1600 & $\tfrac{1}{45}$ & 0 & \texttt{113722}$_{240}$ & 10 & $\tfrac{83}{144}$ & $(0,0,\tfrac{1}{2})$\\
\texttt{113724}$_{480}$ & 408 & $-\tfrac{1}{600}$ & $(0,0,\tfrac{1}{2})$ & \texttt{113725}$_{480}$ & 250 & $\tfrac{247}{720}$ & $(0,0,-\tfrac{1}{2})$\\
\texttt{113744}$_{240}$ & 818 & $\tfrac{1009}{3600}$ & $(0,-1,-\tfrac{1}{2})$ & \texttt{113747}$_{480}$ & 1063 & $-\tfrac{377}{7200}$ & $(0,0,\tfrac{1}{2})$\\
\texttt{113755}$_{240}$ & 534 & $\tfrac{79}{1200}$ & $(0,0,\tfrac{1}{2})$ & \texttt{113757}$_{480}$ & 978 & $\tfrac{173}{1200}$ & $(0,0,-\tfrac{1}{2})$\\
\texttt{113777}$_{240}$ & 1196 & $\tfrac{29}{1800}$ & 0 & \texttt{114414}$_{240}$ & 120 & $\tfrac{2}{5}$ & 0\\
\texttt{114415}$_{240}$ & 34 & $\tfrac{2117}{3600}$ & 0 & \texttt{114416}$_{240}$ & 494 & $\tfrac{187}{3600}$ & 0\\
\texttt{114417}$_{480}$ & 300 & $\tfrac{5}{24}$ & 0 & \texttt{114423}$_{480}$ & 55 & $\tfrac{551}{1440}$ & 0\\
\texttt{114424}$_{480}$ & 1 & $\tfrac{6061}{7200}$ & 0 & \texttt{114427}$_{480}$ & 106 & $\tfrac{983}{3600}$ & $(0,0,-\tfrac{1}{2})$\\
\texttt{114434}$_{480}$ & 409 & $\tfrac{1309}{7200}$ & 0 & \texttt{114437}$_{480}$ & 566 & $\tfrac{673}{3600}$ & $(0,0,-\tfrac{1}{2})$\\
\texttt{114445}$_{480}$ & 191 & $\tfrac{2111}{7200}$ & $(\tfrac{1}{2},0,0)$ & \texttt{114446}$_{480}$ & 599 & $-\tfrac{121}{7200}$ & $(\tfrac{1}{2},0,0)$\\
\texttt{114447}$_{480}$ & 480 & $\tfrac{1}{10}$ & $(0,-1,\tfrac{1}{2})$ & \texttt{114456}$_{480}$ & 545 & $\tfrac{97}{1440}$ & 0\\
\texttt{114513}$_{480}$ & 69 & $\tfrac{1043}{2400}$ & 0 & \texttt{114514}$_{480}$ & 9 & $\tfrac{1703}{2400}$ & 0\\
\texttt{114516}$_{480}$ & 229 & $\tfrac{649}{7200}$ & 0 & \texttt{114527}$_{480}$ & 11 & $\tfrac{3731}{7200}$ & $(0,0,\tfrac{1}{2})$\\
\texttt{114544}$_{240}$ & 120 & $\tfrac{3}{8}$ & $(-\tfrac{1}{2},0,0)$ & \texttt{114547}$_{480}$ & 231 & $-\tfrac{1}{800}$ & $(0,0,\tfrac{1}{2})$\\
\texttt{114557}$_{480}$ & 171 & $\tfrac{817}{2400}$ & $(0,0,-\tfrac{1}{2})$ & \texttt{114612}$_{480}$ & 99 & $\tfrac{493}{2400}$ & 0\\
\texttt{114614}$_{480}$ & 561 & $\tfrac{167}{2400}$ & 0 & \texttt{114615}$_{480}$ & 419 & $\tfrac{1079}{7200}$ & 0\\
\texttt{114637}$_{480}$ & 1261 & $-\tfrac{419}{7200}$ & $(0,0,\tfrac{1}{2})$ & \texttt{114644}$_{240}$ & 840 & $\tfrac{1}{8}$ & $(-\tfrac{1}{2},0,0)$\\
\texttt{114647}$_{480}$ & 1119 & $\tfrac{111}{800}$ & $(0,0,-\tfrac{1}{2})$ & \texttt{114667}$_{480}$ & 1581 & $\tfrac{287}{2400}$ & $(0,0,-\tfrac{1}{2})$\\
\texttt{114711}$_{240}$ & 150 & $\tfrac{17}{48}$ & 0 & \texttt{114712}$_{480}$ & 7 & $\tfrac{5767}{7200}$ & $(0,0,-\tfrac{1}{2})$\\
\texttt{114713}$_{480}$ & 497 & $\tfrac{1577}{7200}$ & $(0,0,-\tfrac{1}{2})$ & \texttt{114714}$_{480}$ & 360 & $\tfrac{1}{5}$ & 0\\
\texttt{114724}$_{480}$ & 118 & $\tfrac{689}{3600}$ & $(0,0,-\tfrac{1}{2})$ & \texttt{114725}$_{480}$ & 48 & $\tfrac{63}{200}$ & $(0,0,\tfrac{1}{2})$\\
\texttt{114734}$_{480}$ & 742 & $\tfrac{41}{3600}$ & $(0,0,\tfrac{1}{2})$ & \texttt{114736}$_{480}$ & 1152 & $\tfrac{27}{200}$ & $(0,0,-\tfrac{1}{2})$\\
\texttt{114745}$_{480}$ & 458 & $\tfrac{1039}{3600}$ & $(0,0,-\tfrac{1}{2})$ & \texttt{114746}$_{480}$ & 1082 & $-\tfrac{329}{3600}$ & $(0,0,\tfrac{1}{2})$\\
\texttt{114747}$_{480}$ & 840 & $\tfrac{3}{10}$ & $(0,0,-1)$ & \texttt{114757}$_{480}$ & 703 & $-\tfrac{917}{7200}$ & $(0,0,1)$\\
\texttt{114767}$_{480}$ & 1193 & $-\tfrac{1387}{7200}$ & $(0,0,1)$ & \texttt{114777}$_{480}$ & 1050 & $\tfrac{1}{16}$ & 0\\
\texttt{115434}$_{480}$ & 267 & $\tfrac{629}{2400}$ & 0 & \texttt{115436}$_{480}$ & 583 & $\tfrac{403}{7200}$ & 0\\
\texttt{115446}$_{480}$ & 485 & $\tfrac{217}{1440}$ & $(-\tfrac{1}{2},0,0)$ & \texttt{115516}$_{240}$ & 104 & $\tfrac{253}{900}$ & 0\\
\texttt{115533}$_{240}$ & 188 & $\tfrac{377}{1800}$ & 0 & \texttt{115544}$_{240}$ & 40 & $\tfrac{179}{360}$ & $(\tfrac{1}{2},0,0)$\\
\texttt{115611}$_{120}$ & 80 & $\tfrac{31}{90}$ & 0 & \texttt{115634}$_{480}$ & 714 & $\tfrac{33}{400}$ & 0\\
\texttt{115666}$_{240}$ & 1508 & $\tfrac{47}{1800}$ & 0 & \texttt{115711}$_{120}$ & 4 & $\tfrac{991}{1800}$ & $(0,0,1)$\\
\texttt{115713}$_{480}$ & 222 & $\tfrac{9}{400}$ & $(0,0,\tfrac{1}{2})$ & \texttt{115714}$_{480}$ & 137 & $\tfrac{2117}{7200}$ & 0\\
\texttt{115733}$_{240}$ & 666 & $\tfrac{281}{1200}$ & $(0,0,-\tfrac{1}{2})$ & \texttt{115736}$_{480}$ & 950 & $-\tfrac{31}{720}$ & $(0,0,\tfrac{1}{2})$\\
\texttt{115744}$_{240}$ & 382 & $\tfrac{611}{3600}$ & $(0,-1,\tfrac{1}{2})$ & \texttt{115746}$_{480}$ & 792 & $\tfrac{91}{600}$ & $(0,0,-\tfrac{1}{2})$\\
\texttt{115766}$_{240}$ & 1190 & $\tfrac{89}{720}$ & $(0,0,-\tfrac{1}{2})$ & \texttt{116424}$_{480}$ & 483 & $\tfrac{101}{2400}$ & 0\\
\texttt{116425}$_{480}$ & 247 & $\tfrac{1267}{7200}$ & 0 & \texttt{116445}$_{480}$ & 595 & $-\tfrac{1}{1440}$ & $(\tfrac{1}{2},0,0)$\\
\texttt{116524}$_{480}$ & 174 & $\tfrac{63}{400}$ & 0 & \texttt{116555}$_{240}$ & 212 & $\tfrac{263}{1800}$ & 0\\
\texttt{116615}$_{240}$ & 1256 & $\tfrac{37}{900}$ & 0 & \texttt{116622}$_{240}$ & 628 & $\tfrac{127}{1800}$ & 0\\
\texttt{116644}$_{240}$ & 1640 & $\tfrac{19}{360}$ & $(-\tfrac{1}{2},0,0)$ & \texttt{116711}$_{120}$ & 764 & $-\tfrac{199}{1800}$ & $(0,0,1)$\\
\texttt{116712}$_{480}$ & 558 & $\tfrac{81}{400}$ & $(0,0,-\tfrac{1}{2})$ & \texttt{116714}$_{480}$ & 1087 & $\tfrac{187}{7200}$ & 0\\
\texttt{116722}$_{240}$ & 234 & $-\tfrac{31}{1200}$ & $(0,0,\tfrac{1}{2})$ & \texttt{116725}$_{480}$ & 730 & $-\tfrac{41}{720}$ & $(0,0,\tfrac{1}{2})$\\
\texttt{116744}$_{240}$ & 1198 & $\tfrac{179}{3600}$ & $(0,1,-\tfrac{1}{2})$ & \texttt{116745}$_{480}$ & 1128 & $-\tfrac{41}{600}$ & $(0,0,\tfrac{1}{2})$\\
\texttt{116755}$_{240}$ & 970 & $\tfrac{127}{720}$ & $(0,0,-\tfrac{1}{2})$ & \texttt{117423}$_{480}$ & 395 & $\tfrac{163}{1440}$ & 0\\
\texttt{117424}$_{480}$ & 218 & $\tfrac{559}{3600}$ & $(0,0,\tfrac{1}{2})$ & \texttt{117434}$_{480}$ & 598 & $-\tfrac{271}{3600}$ & $(0,0,\tfrac{1}{2})$\\
\texttt{117444}$_{480}$ & 525 & $\tfrac{11}{96}$ & 0 & \texttt{117524}$_{480}$ & 35 & $\tfrac{179}{288}$ & $(0,0,-\tfrac{1}{2})$\\
\texttt{117544}$_{240}$ & 238 & $\tfrac{1109}{3600}$ & $(-\tfrac{1}{2},0,-1)$ & \texttt{117634}$_{480}$ & 1475 & $-\tfrac{113}{1440}$ & $(0,0,\tfrac{1}{2})$\\
\texttt{117644}$_{240}$ & 1298 & $\tfrac{799}{3600}$ & $(\tfrac{1}{2},0,-1)$ & \texttt{117711}$_{60}$ & 480 & $\tfrac{3}{5}$ & $(0,0,-2)$\\
\texttt{117712}$_{240}$ & 284 & $\tfrac{41}{1800}$ & $(0,0,1)$ & \texttt{117713}$_{240}$ & 964 & $-\tfrac{329}{1800}$ & $(0,0,1)$\\
\texttt{117714}$_{480}$ & 750 & $\tfrac{5}{48}$ & 0 & \texttt{117722}$_{240}$ & 78 & $\tfrac{201}{400}$ & $(0,0,-\tfrac{1}{2})$\\
\texttt{117724}$_{480}$ & 607 & $\tfrac{1867}{7200}$ & $(0,0,-1)$ & \texttt{117733}$_{240}$ & 1182 & $-\tfrac{31}{400}$ & $(0,0,\tfrac{1}{2})$\\
\texttt{117734}$_{480}$ & 1097 & $\tfrac{1877}{7200}$ & $(0,0,-1)$ & \texttt{117744}$_{240}$ & 960 & $-\tfrac{3}{10}$ & $(0,1,\tfrac{3}{2})$\\
\texttt{122337}$_{240}$ & 60 & $\tfrac{17}{40}$ & 0 & \texttt{122346}$_{480}$ & 82 & $\tfrac{941}{3600}$ & 0\\
\texttt{122367}$_{480}$ & 208 & $\tfrac{29}{225}$ & 0 & \texttt{122376}$_{480}$ & 308 & $\tfrac{287}{1800}$ & 0\\
\texttt{122446}$_{240}$ & 8 & $\tfrac{1067}{1800}$ & $(\tfrac{1}{2},0,0)$ & \texttt{122467}$_{480}$ & 98 & $\tfrac{589}{3600}$ & 0\\
\texttt{122477}$_{480}$ & 21 & $\tfrac{1267}{2400}$ & 0 & \texttt{122666}$_{240}$ & 704 & $\tfrac{7}{225}$ & 0\\
\texttt{122676}$_{480}$ & 572 & $\tfrac{113}{1800}$ & 0 & \texttt{122777}$_{240}$ & 162 & $\tfrac{397}{1200}$ & $(0,0,-\tfrac{1}{2})$\\
\texttt{123324}$_{240}$ & 190 & $\tfrac{163}{720}$ & 0 & \texttt{123344}$_{480}$ & 394 & $\tfrac{497}{3600}$ & 0\\
\texttt{123351}$_{240}$ & 112 & $\tfrac{86}{225}$ & 0 & \texttt{123371}$_{240}$ & 344 & $\tfrac{103}{900}$ & 0\\
\texttt{123424}$_{480}$ & 14 & $\tfrac{2167}{3600}$ & 0 & \texttt{123441}$_{480}$ & 30 & $\tfrac{121}{240}$ & 0\\
\texttt{123454}$_{480}$ & 94 & $\tfrac{737}{3600}$ & $(\tfrac{1}{2},0,0)$ & \texttt{123471}$_{480}$ & 109 & $\tfrac{349}{7200}$ & $(0,0,\tfrac{1}{2})$\\
\texttt{123621}$_{240}$ & 272 & $\tfrac{31}{225}$ & 0 & \texttt{123644}$_{480}$ & 886 & $\tfrac{53}{3600}$ & $(\tfrac{1}{2},0,0)$\\
\texttt{123651}$_{480}$ & 420 & $\tfrac{13}{120}$ & 0 & \texttt{123674}$_{480}$ & 1026 & $\tfrac{31}{1200}$ & 0\\
\texttt{123721}$_{240}$ & 56 & $\tfrac{427}{900}$ & 0 & \texttt{123741}$_{480}$ & 299 & $\tfrac{659}{7200}$ & $(0,0,\tfrac{1}{2})$\\
\texttt{123754}$_{480}$ & 486 & $\tfrac{161}{1200}$ & 0 & \texttt{123774}$_{480}$ & 685 & $\tfrac{77}{1440}$ & 0\\
\texttt{124335}$_{480}$ & 182 & $\tfrac{931}{3600}$ & 0 & \texttt{124344}$_{480}$ & 240 & $\tfrac{1}{4}$ & 0\\
\texttt{124347}$_{480}$ & 450 & $\tfrac{19}{240}$ & 0 & \texttt{124353}$_{480}$ & 298 & $\tfrac{689}{3600}$ & 0\\
\texttt{124364}$_{480}$ & 466 & $\tfrac{173}{3600}$ & 0 & \texttt{124373}$_{480}$ & 463 & $-\tfrac{257}{7200}$ & $(0,0,\tfrac{1}{2})$\\
\texttt{124374}$_{480}$ & 365 & $\tfrac{361}{1440}$ & $(0,0,-\tfrac{1}{2})$ & \texttt{124427}$_{240}$ & 44 & $\tfrac{551}{1800}$ & 0\\
\texttt{124434}$_{480}$ & 100 & $\tfrac{77}{360}$ & 0 & \texttt{124443}$_{480}$ & 116 & $\tfrac{209}{1800}$ & 0\\
\texttt{124444}$_{480}$ & 60 & $\tfrac{3}{8}$ & 0 & \texttt{124445}$_{480}$ & 4 & $\tfrac{319}{450}$ & $(\tfrac{1}{2},0,0)$\\
\texttt{124454}$_{480}$ & 20 & $\tfrac{11}{18}$ & $(-\tfrac{1}{2},0,0)$ & \texttt{124461}$_{240}$ & 76 & $\tfrac{319}{1800}$ & 0\\
\texttt{124471}$_{240}$ & 14 & $\tfrac{2767}{3600}$ & $(0,0,-1)$ & \texttt{124474}$_{480}$ & 119 & $\tfrac{1739}{7200}$ & $(0,-1,-\tfrac{1}{2})$\\
\texttt{124475}$_{480}$ & 65 & $\tfrac{65}{288}$ & $(0,0,\tfrac{1}{2})$ & \texttt{124621}$_{240}$ & 22 & $\tfrac{1691}{3600}$ & 0\\
\texttt{124634}$_{480}$ & 674 & $\tfrac{217}{3600}$ & 0 & \texttt{124643}$_{480}$ & 848 & $\tfrac{31}{900}$ & 0\\
\texttt{124644}$_{480}$ & 600 & $\tfrac{1}{8}$ & $(-\tfrac{1}{2},0,0)$ & \texttt{124645}$_{480}$ & 352 & $\tfrac{59}{900}$ & 0\\
\texttt{124654}$_{480}$ & 526 & $\tfrac{323}{3600}$ & 0 & \texttt{124667}$_{480}$ & 1178 & $\tfrac{49}{3600}$ & 0\\
\texttt{124673}$_{480}$ & 953 & $\tfrac{1073}{7200}$ & $(0,0,-\tfrac{1}{2})$ & \texttt{124674}$_{480}$ & 717 & $-\tfrac{67}{800}$ & $(0,0,\tfrac{1}{2})$\\
\texttt{124677}$_{480}$ & 1101 & $-\tfrac{193}{2400}$ & $(0,0,\tfrac{1}{2})$ & \texttt{124724}$_{480}$ & 80 & $\tfrac{47}{180}$ & 0\\
\texttt{124733}$_{480}$ & 594 & $\tfrac{119}{1200}$ & 0 & \texttt{124741}$_{480}$ & 81 & $\tfrac{947}{2400}$ & $(0,0,-\tfrac{1}{2})$\\
\texttt{124744}$_{480}$ & 360 & $\tfrac{1}{8}$ & $(0,1,0)$ & \texttt{124747}$_{480}$ & 639 & $\tfrac{51}{800}$ & 0\\
\texttt{124755}$_{480}$ & 126 & $\tfrac{421}{1200}$ & 0 & \texttt{124764}$_{480}$ & 640 & $\tfrac{7}{180}$ & 0\\
\texttt{124774}$_{480}$ & 502 & $\tfrac{41}{3600}$ & $(0,0,\tfrac{1}{2})$ & \texttt{124775}$_{480}$ & 325 & $\tfrac{61}{288}$ & $(0,0,-\tfrac{1}{2})$\\
\texttt{124777}$_{480}$ & 713 & $\tfrac{1733}{7200}$ & $(0,0,-1)$ & \texttt{125337}$_{480}$ & 368 & $\tfrac{34}{225}$ & 0\\
\texttt{125344}$_{480}$ & 86 & $\tfrac{1483}{3600}$ & 0 & \texttt{125364}$_{480}$ & 290 & $\tfrac{89}{720}$ & 0\\
\texttt{125434}$_{480}$ & 26 & $\tfrac{1873}{3600}$ & 0 & \texttt{125447}$_{480}$ & 90 & $\tfrac{3}{16}$ & $(0,1,0)$\\
\texttt{125464}$_{480}$ & 106 & $\tfrac{443}{3600}$ & $(\tfrac{1}{2},0,0)$ & \texttt{125637}$_{480}$ & 780 & $\tfrac{7}{120}$ & 0\\
\texttt{125644}$_{480}$ & 314 & $\tfrac{307}{3600}$ & $(\tfrac{1}{2},0,0)$ & \texttt{125667}$_{480}$ & 928 & $\tfrac{13}{450}$ & 0\\
\texttt{125734}$_{480}$ & 234 & $\tfrac{259}{1200}$ & 0 & \texttt{125747}$_{480}$ & 421 & $\tfrac{1081}{7200}$ & 0\\
\texttt{125767}$_{480}$ & 664 & $\tfrac{53}{900}$ & 0 & \texttt{126342}$_{480}$ & 398 & $\tfrac{319}{3600}$ & 0\\
\texttt{126442}$_{480}$ & 112 & $\tfrac{119}{900}$ & 0 & \texttt{126452}$_{480}$ & 38 & $\tfrac{1579}{3600}$ & 0\\
\texttt{126622}$_{240}$ & 496 & $\tfrac{31}{450}$ & 0 & \texttt{126641}$_{240}$ & 1102 & $\tfrac{71}{3600}$ & 0\\
\texttt{126651}$_{240}$ & 992 & $\tfrac{17}{450}$ & 0 & \texttt{126722}$_{480}$ & 148 & $\tfrac{427}{1800}$ & 0\\
\texttt{126741}$_{480}$ & 699 & $\tfrac{353}{2400}$ & $(0,0,-\tfrac{1}{2})$ & \texttt{126752}$_{480}$ & 412 & $\tfrac{253}{1800}$ & 0\\
\texttt{127335}$_{480}$ & 452 & $\tfrac{143}{1800}$ & 0 & \texttt{127342}$_{480}$ & 127 & $\tfrac{2767}{7200}$ & $(0,0,-\tfrac{1}{2})$\\
\texttt{127344}$_{480}$ & 475 & $\tfrac{43}{288}$ & $(0,0,-\tfrac{1}{2})$ & \texttt{127432}$_{480}$ & 54 & $\tfrac{409}{1200}$ & 0\\
\texttt{127442}$_{480}$ & 3 & $\tfrac{1561}{2400}$ & $(0,0,\tfrac{1}{2})$ & \texttt{127445}$_{480}$ & 115 & $\tfrac{83}{1440}$ & $(-\tfrac{1}{2},0,\tfrac{1}{2})$\\
\texttt{127634}$_{480}$ & 1134 & $\tfrac{29}{1200}$ & 0 & \texttt{127644}$_{480}$ & 1083 & $\tfrac{127}{800}$ & $(-\tfrac{1}{2},0,-\tfrac{1}{2})$\\
\texttt{127645}$_{480}$ & 847 & $-\tfrac{353}{7200}$ & $(0,0,\tfrac{1}{2})$ & \texttt{127725}$_{240}$ & 296 & $\tfrac{157}{900}$ & 0\\
\texttt{127731}$_{240}$ & 560 & $\tfrac{7}{90}$ & 0 & \texttt{127742}$_{480}$ & 219 & $\tfrac{11}{800}$ & $(0,0,\tfrac{1}{2})$\\
\texttt{127744}$_{480}$ & 681 & $\tfrac{47}{2400}$ & $(0,-1,\tfrac{1}{2})$ & \texttt{127745}$_{480}$ & 539 & $\tfrac{1619}{7200}$ & $(0,0,-\tfrac{1}{2})$\\
\texttt{133445}$_{240}$ & 872 & $\tfrac{113}{1800}$ & $(\tfrac{1}{2},0,0)$ & \texttt{133457}$_{480}$ & 962 & $\tfrac{301}{3600}$ & 0\\
\texttt{133477}$_{480}$ & 1491 & $\tfrac{77}{2400}$ & 0 & \texttt{133555}$_{240}$ & 416 & $\tfrac{43}{225}$ & 0\\
\texttt{133575}$_{480}$ & 1012 & $\tfrac{103}{1800}$ & 0 & \texttt{133777}$_{240}$ & 1938 & $-\tfrac{107}{1200}$ & $(0,0,\tfrac{1}{2})$\\
\texttt{134437}$_{240}$ & 1196 & $\tfrac{119}{1800}$ & 0 & \texttt{134442}$_{480}$ & 436 & $\tfrac{259}{1800}$ & 0\\
\texttt{134444}$_{480}$ & 780 & $\tfrac{1}{8}$ & 0 & \texttt{134446}$_{480}$ & 1124 & $\tfrac{73}{900}$ & $(-\tfrac{1}{2},0,0)$\\
\texttt{134451}$_{240}$ & 364 & $\tfrac{391}{1800}$ & 0 & \texttt{134464}$_{480}$ & 1460 & $\tfrac{11}{180}$ & $(-\tfrac{1}{2},0,0)$\\
\texttt{134471}$_{240}$ & 994 & $\tfrac{977}{3600}$ & $(0,0,-1)$ & \texttt{134474}$_{480}$ & 1151 & $-\tfrac{709}{7200}$ & $(0,1,\tfrac{1}{2})$\\
\texttt{134476}$_{480}$ & 1505 & $\tfrac{181}{1440}$ & $(0,0,-\tfrac{1}{2})$ & \texttt{134531}$_{240}$ & 598 & $\tfrac{539}{3600}$ & 0\\
\texttt{134542}$_{480}$ & 208 & $\tfrac{221}{900}$ & 0 & \texttt{134544}$_{480}$ & 600 & $\tfrac{1}{8}$ & $(\tfrac{1}{2},0,0)$\\
\texttt{134546}$_{480}$ & 992 & $\tfrac{49}{900}$ & 0 & \texttt{134557}$_{480}$ & 602 & $\tfrac{481}{3600}$ & 0\\
\texttt{134564}$_{480}$ & 1174 & $\tfrac{107}{3600}$ & 0 & \texttt{134572}$_{480}$ & 617 & $-\tfrac{223}{7200}$ & $(0,0,\tfrac{1}{2})$\\
\texttt{134574}$_{480}$ & 933 & $\tfrac{157}{800}$ & $(0,0,-\tfrac{1}{2})$ & \texttt{134577}$_{480}$ & 1131 & $-\tfrac{143}{2400}$ & $(0,0,\tfrac{1}{2})$\\
\texttt{134734}$_{480}$ & 1360 & $\tfrac{7}{180}$ & 0 & \texttt{134741}$_{480}$ & 969 & $-\tfrac{157}{2400}$ & $(0,0,\tfrac{1}{2})$\\
\texttt{134744}$_{480}$ & 1080 & $\tfrac{1}{8}$ & $(0,-1,0)$ & \texttt{134747}$_{480}$ & 1191 & $\tfrac{19}{800}$ & 0\\
\texttt{134754}$_{480}$ & 800 & $\tfrac{11}{180}$ & 0 & \texttt{134766}$_{480}$ & 2106 & $\tfrac{11}{1200}$ & 0\\
\texttt{134774}$_{480}$ & 1562 & $-\tfrac{329}{3600}$ & $(0,0,\tfrac{1}{2})$ & \texttt{134776}$_{480}$ & 1765 & $-\tfrac{127}{1440}$ & $(0,0,\tfrac{1}{2})$\\
\texttt{134777}$_{480}$ & 1663 & $\tfrac{1723}{7200}$ & $(0,0,-1)$ & \texttt{135443}$_{480}$ & 688 & $\tfrac{101}{900}$ & 0\\
\texttt{135463}$_{480}$ & 1378 & $\tfrac{149}{3600}$ & 0 & \texttt{135533}$_{240}$ & 784 & $\tfrac{49}{450}$ & 0\\
\texttt{135541}$_{240}$ & 238 & $\tfrac{1079}{3600}$ & 0 & \texttt{135561}$_{240}$ & 832 & $\tfrac{37}{450}$ & 0\\
\texttt{135733}$_{480}$ & 1148 & $\tfrac{77}{1800}$ & 0 & \texttt{135741}$_{480}$ & 669 & $\tfrac{463}{2400}$ & $(0,0,-\tfrac{1}{2})$\\
\texttt{135763}$_{480}$ & 1708 & $\tfrac{37}{1800}$ & 0 & \texttt{136447}$_{480}$ & 1530 & $\tfrac{7}{80}$ & $(0,-1,0)$\\
\texttt{136454}$_{480}$ & 1166 & $\tfrac{313}{3600}$ & $(-\tfrac{1}{2},0,0)$ & \texttt{136544}$_{480}$ & 1106 & $-\tfrac{17}{3600}$ & $(\tfrac{1}{2},0,0)$\\
\texttt{136557}$_{480}$ & 1088 & $\tfrac{23}{450}$ & 0 & \texttt{136747}$_{480}$ & 2051 & $\tfrac{71}{7200}$ & 0\\
\texttt{136757}$_{480}$ & 1816 & $\tfrac{17}{900}$ & 0 & \texttt{137443}$_{480}$ & 1227 & $\tfrac{409}{2400}$ & $(0,0,-\tfrac{1}{2})$\\
\texttt{137446}$_{480}$ & 1445 & $-\tfrac{31}{288}$ & $(\tfrac{1}{2},0,\tfrac{1}{2})$ & \texttt{137544}$_{480}$ & 867 & $\tfrac{23}{800}$ & $(-\tfrac{1}{2},0,\tfrac{1}{2})$\\
\texttt{137546}$_{480}$ & 1183 & $\tfrac{943}{7200}$ & $(0,0,-\tfrac{1}{2})$ & \texttt{137736}$_{240}$ & 2024 & $\tfrac{13}{900}$ & 0\\
\texttt{137743}$_{480}$ & 1629 & $-\tfrac{59}{800}$ & $(0,0,\tfrac{1}{2})$ & \texttt{137744}$_{480}$ & 1569 & $\tfrac{143}{2400}$ & $(0,1,-\tfrac{1}{2})$\\
\texttt{137746}$_{480}$ & 1789 & $\tfrac{829}{7200}$ & $(0,0,-\tfrac{1}{2})$ & \texttt{144441}$_{120}$ & 400 & $\tfrac{2}{9}$ & 0\\
\texttt{144447}$_{240}$ & 800 & $\tfrac{1}{9}$ & 0 & \texttt{144452}$_{480}$ & 76 & $\tfrac{799}{1800}$ & 0\\
\texttt{144454}$_{480}$ & 420 & $\tfrac{1}{4}$ & $(-\tfrac{1}{2},0,0)$ & \texttt{144456}$_{480}$ & 764 & $\tfrac{7}{225}$ & $(\tfrac{1}{2},0,0)$\\
\texttt{144463}$_{480}$ & 1196 & $\tfrac{29}{1800}$ & 0 & \texttt{144464}$_{480}$ & 1140 & $0$ & $(\tfrac{1}{2},0,0)$\\
\texttt{144465}$_{480}$ & 1084 & $\tfrac{53}{900}$ & $(-\tfrac{1}{2},0,0)$ & \texttt{144471}$_{240}$ & 720 & $-\tfrac{1}{10}$ & $(0,0,1)$\\
\texttt{144472}$_{480}$ & 601 & $\tfrac{961}{7200}$ & $(0,0,-\tfrac{1}{2})$ & \texttt{144473}$_{480}$ & 1009 & $-\tfrac{191}{7200}$ & $(0,0,\tfrac{1}{2})$\\
\texttt{144475}$_{480}$ & 791 & $\tfrac{1811}{7200}$ & $(-\tfrac{1}{2},0,-\tfrac{1}{2})$ & \texttt{144476}$_{480}$ & 1199 & $-\tfrac{421}{7200}$ & $(-\tfrac{1}{2},0,\tfrac{1}{2})$\\
\texttt{144477}$_{480}$ & 1080 & $\tfrac{1}{10}$ & $(0,1,-\tfrac{1}{2})$ & \texttt{144555}$_{240}$ & 152 & $\tfrac{743}{1800}$ & $(-\tfrac{1}{2},0,0)$\\
\texttt{144561}$_{480}$ & 750 & $\tfrac{13}{240}$ & 0 & \texttt{144572}$_{480}$ & 355 & $\tfrac{407}{1440}$ & $(0,0,-\tfrac{1}{2})$\\
\texttt{144575}$_{480}$ & 573 & $\tfrac{11}{2400}$ & $(\tfrac{1}{2},0,\tfrac{1}{2})$ & \texttt{144577}$_{480}$ & 831 & $\tfrac{497}{2400}$ & $(0,-1,-\tfrac{1}{2})$\\
\texttt{144666}$_{240}$ & 2168 & $\tfrac{77}{1800}$ & $(-\tfrac{1}{2},0,0)$ & \texttt{144673}$_{480}$ & 1685 & $\tfrac{169}{1440}$ & $(0,0,-\tfrac{1}{2})$\\
\texttt{144676}$_{480}$ & 1797 & $\tfrac{179}{2400}$ & $(\tfrac{1}{2},0,-\tfrac{1}{2})$ & \texttt{144677}$_{480}$ & 1719 & $-\tfrac{367}{2400}$ & $(0,1,\tfrac{1}{2})$\\
\texttt{144774}$_{480}$ & 1275 & $\tfrac{5}{96}$ & 0 & \texttt{144775}$_{240}$ & 1202 & $\tfrac{691}{3600}$ & $(\tfrac{1}{2},0,-1)$\\
\texttt{144776}$_{240}$ & 1582 & $\tfrac{1001}{3600}$ & $(-\tfrac{1}{2},0,-1)$ & \texttt{144777}$_{240}$ & 1440 & $-\tfrac{1}{5}$ & $(0,-1,\tfrac{3}{2})$\\
\texttt{145464}$_{480}$ & 960 & $0$ & $(1,0,0)$ & \texttt{145467}$_{480}$ & 1170 & $\tfrac{7}{240}$ & 0\\
\texttt{145547}$_{240}$ & 476 & $\tfrac{299}{1800}$ & 0 & \texttt{145564}$_{480}$ & 754 & $\tfrac{407}{3600}$ & $(-\tfrac{1}{2},0,0)$\\
\texttt{145644}$_{480}$ & 820 & $\tfrac{41}{360}$ & $(-1,0,0)$ & \texttt{145666}$_{480}$ & 1882 & $\tfrac{41}{3600}$ & 0\\
\texttt{145742}$_{480}$ & 295 & $\tfrac{7}{288}$ & $(0,0,\tfrac{1}{2})$ & \texttt{145744}$_{480}$ & 649 & $\tfrac{889}{7200}$ & $(\tfrac{1}{2},1,-\tfrac{1}{2})$\\
\texttt{145747}$_{480}$ & 806 & $-\tfrac{167}{3600}$ & $(0,0,\tfrac{1}{2})$ & \texttt{145764}$_{480}$ & 1325 & $-\tfrac{179}{1440}$ & $(\tfrac{1}{2},0,\tfrac{1}{2})$\\
\texttt{145766}$_{480}$ & 1673 & $-\tfrac{607}{7200}$ & $(0,0,\tfrac{1}{2})$ & \texttt{145767}$_{480}$ & 1501 & $\tfrac{1021}{7200}$ & $(0,0,-\tfrac{1}{2})$\\
\texttt{146555}$_{480}$ & 542 & $\tfrac{391}{3600}$ & 0 & \texttt{146647}$_{240}$ & 2204 & $\tfrac{11}{1800}$ & 0\\
\texttt{146654}$_{480}$ & 1814 & $-\tfrac{83}{3600}$ & $(\tfrac{1}{2},0,0)$ & \texttt{146743}$_{480}$ & 1735 & $\tfrac{179}{1440}$ & $(0,0,-\tfrac{1}{2})$\\
\texttt{146744}$_{480}$ & 1681 & $-\tfrac{479}{7200}$ & $(\tfrac{1}{2},-1,\tfrac{1}{2})$ & \texttt{146747}$_{480}$ & 1786 & $-\tfrac{337}{3600}$ & $(0,0,\tfrac{1}{2})$\\
\texttt{146754}$_{480}$ & 1435 & $-\tfrac{37}{1440}$ & $(-\tfrac{1}{2},0,\tfrac{1}{2})$ & \texttt{146755}$_{480}$ & 1337 & $\tfrac{977}{7200}$ & $(0,0,-\tfrac{1}{2})$\\
\texttt{146757}$_{480}$ & 1691 & $\tfrac{851}{7200}$ & $(0,0,-\tfrac{1}{2})$ & \texttt{147744}$_{240}$ & 1320 & $\tfrac{2}{5}$ & $(0,-2,-1)$\\
\texttt{147745}$_{240}$ & 1234 & $-\tfrac{583}{3600}$ & $(0,0,1)$ & \texttt{147746}$_{240}$ & 1694 & $-\tfrac{713}{3600}$ & $(0,0,1)$\\
\texttt{147747}$_{480}$ & 1500 & $\tfrac{1}{24}$ & 0 & \texttt{155667}$_{240}$ & 1500 & $\tfrac{1}{40}$ & 0\\
\texttt{156654}$_{240}$ & 1630 & $\tfrac{19}{720}$ & 0 & \texttt{222666}$_{80}$ & 456 & $\tfrac{29}{300}$ & 0\\
\texttt{222667}$_{240}$ & 224 & $\tfrac{22}{225}$ & 0 & \texttt{222677}$_{240}$ & 92 & $\tfrac{533}{1800}$ & 0\\
\texttt{222777}$_{80}$ & 6 & $\tfrac{237}{400}$ & $(0,0,\tfrac{1}{2})$ & \texttt{223644}$_{240}$ & 588 & $\tfrac{37}{300}$ & $(-\tfrac{1}{2},0,0)$\\
\texttt{223744}$_{240}$ & 232 & $\tfrac{43}{1800}$ & $(0,1,0)$ & \texttt{224634}$_{480}$ & 492 & $\tfrac{61}{600}$ & 0\\
\texttt{224637}$_{480}$ & 682 & $\tfrac{161}{3600}$ & 0 & \texttt{224644}$_{240}$ & 360 & $\tfrac{1}{8}$ & $(\tfrac{1}{2},0,0)$\\
\texttt{224645}$_{480}$ & 228 & $\tfrac{119}{600}$ & 0 & \texttt{224647}$_{480}$ & 608 & $\tfrac{31}{900}$ & 0\\
\texttt{224734}$_{480}$ & 194 & $\tfrac{637}{3600}$ & 0 & \texttt{224744}$_{240}$ & 120 & $\tfrac{3}{8}$ & $(0,-1,0)$\\
\texttt{224745}$_{480}$ & 46 & $\tfrac{1703}{3600}$ & 0 & \texttt{224747}$_{480}$ & 237 & $\tfrac{139}{2400}$ & 0\\
\texttt{224757}$_{480}$ & 186 & $\tfrac{57}{400}$ & 0 & \texttt{225644}$_{240}$ & 132 & $\tfrac{49}{150}$ & $(-\tfrac{1}{2},0,0)$\\
\texttt{225647}$_{480}$ & 322 & $\tfrac{341}{3600}$ & 0 & \texttt{225677}$_{240}$ & 548 & $\tfrac{167}{1800}$ & 0\\
\texttt{225744}$_{240}$ & 8 & $\tfrac{1127}{1800}$ & $(0,1,0)$ & \texttt{225747}$_{480}$ & 113 & $\tfrac{2093}{7200}$ & 0\\
\texttt{225777}$_{240}$ & 230 & $-\tfrac{7}{720}$ & $(0,0,\tfrac{1}{2})$ & \texttt{226622}$_{60}$ & 264 & $\tfrac{61}{300}$ & 0\\
\texttt{226722}$_{120}$ & 16 & $\tfrac{271}{450}$ & 0 & \texttt{227644}$_{240}$ & 712 & $-\tfrac{17}{1800}$ & $(\tfrac{1}{2},0,0)$\\
\texttt{227645}$_{480}$ & 638 & $\tfrac{199}{3600}$ & 0 & \texttt{227724}$_{240}$ & 142 & $\tfrac{551}{3600}$ & 0\\
\texttt{227725}$_{240}$ & 32 & $\tfrac{227}{450}$ & 0 & \texttt{227755}$_{240}$ & 180 & $\tfrac{5}{24}$ & 0\\
\texttt{233445}$_{240}$ & 732 & $\tfrac{7}{75}$ & $(-\tfrac{1}{2},0,0)$ & \texttt{233447}$_{240}$ & 872 & $-\tfrac{7}{1800}$ & $(0,1,0)$\\
\texttt{234434}$_{240}$ & 786 & $\tfrac{91}{1200}$ & 0 & \texttt{234436}$_{240}$ & 972 & $\tfrac{7}{200}$ & 0\\
\texttt{234442}$_{240}$ & 366 & $\tfrac{101}{1200}$ & 0 & \texttt{234444}$_{480}$ & 540 & $\tfrac{1}{8}$ & 0\\
\texttt{234446}$_{480}$ & 714 & $-\tfrac{11}{1200}$ & $(\tfrac{1}{2},0,0)$ & \texttt{234447}$_{480}$ & 644 & $\tfrac{133}{900}$ & $(0,-1,0)$\\
\texttt{234452}$_{240}$ & 108 & $\tfrac{63}{200}$ & 0 & \texttt{234454}$_{480}$ & 294 & $\tfrac{239}{1200}$ & $(-\tfrac{1}{2},0,0)$\\
\texttt{234474}$_{480}$ & 964 & $\tfrac{83}{900}$ & $(0,-1,0)$ & \texttt{234477}$_{480}$ & 1025 & $\tfrac{49}{1440}$ & 0\\
\texttt{235443}$_{480}$ & 348 & $\tfrac{49}{600}$ & 0 & \texttt{235473}$_{480}$ & 782 & $\tfrac{271}{3600}$ & 0\\
\texttt{237443}$_{480}$ & 986 & $\tfrac{133}{3600}$ & 0 & \texttt{237446}$_{480}$ & 1066 & $\tfrac{203}{3600}$ & $(-\tfrac{1}{2},0,0)$\\
\texttt{237454}$_{480}$ & 686 & $\tfrac{193}{3600}$ & $(\tfrac{1}{2},0,0)$ & \texttt{237456}$_{480}$ & 890 & $\tfrac{29}{720}$ & 0\\
\texttt{244442}$_{120}$ & 144 & $\tfrac{8}{25}$ & 0 & \texttt{244443}$_{240}$ & 432 & $\tfrac{4}{25}$ & 0\\
\texttt{244445}$_{240}$ & 288 & $\tfrac{6}{25}$ & 0 & \texttt{244446}$_{240}$ & 576 & $\tfrac{2}{25}$ & 0\\
\texttt{244452}$_{240}$ & 6 & $\tfrac{881}{1200}$ & 0 & \texttt{244454}$_{480}$ & 180 & $\tfrac{1}{4}$ & $(\tfrac{1}{2},0,0)$\\
\texttt{244456}$_{480}$ & 354 & $\tfrac{109}{1200}$ & $(-\tfrac{1}{2},0,0)$ & \texttt{244457}$_{480}$ & 284 & $\tfrac{22}{225}$ & $(0,1,0)$\\
\texttt{244473}$_{480}$ & 700 & $\tfrac{17}{360}$ & 0 & \texttt{244474}$_{480}$ & 660 & $0$ & $(0,1,0)$\\
\texttt{244475}$_{480}$ & 620 & $\tfrac{1}{36}$ & $(\tfrac{1}{2},0,0)$ & \texttt{244477}$_{480}$ & 719 & $-\tfrac{361}{7200}$ & $(0,1,0)$\\
\texttt{244555}$_{240}$ & 12 & $\tfrac{193}{300}$ & $(\tfrac{1}{2},0,0)$ & \texttt{244557}$_{240}$ & 152 & $\tfrac{623}{1800}$ & $(0,-1,0)$\\
\texttt{244572}$_{480}$ & 254 & $\tfrac{427}{3600}$ & 0 & \texttt{244575}$_{480}$ & 334 & $\tfrac{497}{3600}$ & $(-\tfrac{1}{2},0,0)$\\
\texttt{244776}$_{240}$ & 1240 & $-\tfrac{7}{360}$ & $(\tfrac{1}{2},0,0)$ & \texttt{244777}$_{240}$ & 1202 & $\tfrac{661}{3600}$ & $(0,-1,-\tfrac{1}{2})$\\
\texttt{245474}$_{480}$ & 480 & $\tfrac{1}{4}$ & $(-\tfrac{1}{2},-1,0)$ & \texttt{245476}$_{480}$ & 706 & $\tfrac{113}{3600}$ & 0\\
\texttt{245477}$_{480}$ & 605 & $\tfrac{133}{1440}$ & 0 & \texttt{245544}$_{240}$ & 66 & $\tfrac{631}{1200}$ & $(-1,0,0)$\\
\texttt{245546}$_{240}$ & 252 & $\tfrac{27}{200}$ & 0 & \texttt{245572}$_{240}$ & 70 & $\tfrac{59}{144}$ & 0\\
\texttt{245574}$_{480}$ & 274 & $\tfrac{527}{3600}$ & $(\tfrac{1}{2},0,0)$ & \texttt{245744}$_{480}$ & 356 & $\tfrac{149}{1800}$ & $(\tfrac{1}{2},-1,0)$\\
\texttt{245774}$_{480}$ & 845 & $\tfrac{97}{1440}$ & $(-\tfrac{1}{2},0,0)$ & \texttt{245777}$_{480}$ & 912 & $-\tfrac{13}{200}$ & $(0,0,\tfrac{1}{2})$\\
\texttt{247555}$_{480}$ & 178 & $\tfrac{989}{3600}$ & 0 & \texttt{247557}$_{480}$ & 343 & $\tfrac{643}{7200}$ & 0\\
\texttt{247746}$_{240}$ & 1276 & $\tfrac{19}{1800}$ & 0 & \texttt{247747}$_{240}$ & 1214 & $\tfrac{67}{3600}$ & 0\\
\texttt{247754}$_{480}$ & 870 & $-\tfrac{1}{48}$ & $(0,1,0)$ & \texttt{247757}$_{480}$ & 949 & $\tfrac{169}{7200}$ & 0\\
\texttt{255775}$_{240}$ & 508 & $\tfrac{157}{1800}$ & 0 & \texttt{255777}$_{240}$ & 710 & $\tfrac{137}{720}$ & $(0,0,-\tfrac{1}{2})$\\
\texttt{257755}$_{240}$ & 592 & $\tfrac{26}{225}$ & 0 & \texttt{257757}$_{240}$ & 824 & $\tfrac{43}{900}$ & 0\\
\texttt{333555}$_{80}$ & 696 & $\tfrac{19}{300}$ & 0 & \texttt{333557}$_{240}$ & 1376 & $\tfrac{13}{225}$ & 0\\
\texttt{333577}$_{240}$ & 1972 & $\tfrac{43}{1800}$ & 0 & \texttt{333777}$_{80}$ & 2454 & $\tfrac{53}{400}$ & $(0,0,-\tfrac{1}{2})$\\
\texttt{334544}$_{240}$ & 1080 & $\tfrac{1}{8}$ & $(-\tfrac{1}{2},0,0)$ & \texttt{334546}$_{480}$ & 1812 & $\tfrac{11}{600}$ & 0\\
\texttt{334547}$_{480}$ & 1472 & $\tfrac{49}{900}$ & 0 & \texttt{334744}$_{240}$ & 1560 & $\tfrac{1}{8}$ & $(0,-1,0)$\\
\texttt{334746}$_{480}$ & 2134 & $\tfrac{47}{3600}$ & 0 & \texttt{334747}$_{480}$ & 1893 & $\tfrac{91}{2400}$ & 0\\
\texttt{334767}$_{480}$ & 2526 & $\tfrac{7}{400}$ & 0 & \texttt{335533}$_{60}$ & 1464 & $\tfrac{11}{300}$ & 0\\
\texttt{335733}$_{120}$ & 1744 & $\tfrac{19}{450}$ & 0 & \texttt{336544}$_{240}$ & 1428 & $\tfrac{17}{300}$ & $(-\tfrac{1}{2},0,0)$\\
\texttt{336547}$_{480}$ & 1862 & $\tfrac{91}{3600}$ & 0 & \texttt{336577}$_{240}$ & 2132 & $\tfrac{23}{1800}$ & 0\\
\texttt{336744}$_{240}$ & 2248 & $-\tfrac{83}{1800}$ & $(0,1,0)$ & \texttt{336747}$_{480}$ & 2657 & $\tfrac{77}{7200}$ & 0\\
\texttt{336777}$_{240}$ & 2890 & $-\tfrac{13}{144}$ & $(0,0,\tfrac{1}{2})$ & \texttt{337544}$_{240}$ & 1288 & $\tfrac{37}{1800}$ & $(\tfrac{1}{2},0,0)$\\
\texttt{337546}$_{480}$ & 1978 & $\tfrac{89}{3600}$ & 0 & \texttt{337734}$_{240}$ & 2158 & $\tfrac{119}{3600}$ & 0\\
\texttt{337736}$_{240}$ & 2752 & $\tfrac{7}{450}$ & 0 & \texttt{337766}$_{240}$ & 3060 & $\tfrac{1}{120}$ & 0\\
\texttt{344443}$_{120}$ & 1296 & $\tfrac{2}{25}$ & 0 & \texttt{344445}$_{240}$ & 864 & $\tfrac{3}{25}$ & 0\\
\texttt{344446}$_{240}$ & 1728 & $\tfrac{1}{25}$ & 0 & \texttt{344463}$_{240}$ & 1866 & $\tfrac{31}{1200}$ & 0\\
\texttt{344464}$_{480}$ & 1620 & $0$ & $(\tfrac{1}{2},0,0)$ & \texttt{344465}$_{480}$ & 1374 & $-\tfrac{1}{1200}$ & $(\tfrac{1}{2},0,0)$\\
\texttt{344467}$_{480}$ & 2044 & $-\tfrac{13}{225}$ & $(0,1,0)$ & \texttt{344474}$_{480}$ & 1380 & $0$ & $(0,1,0)$\\
\texttt{344476}$_{480}$ & 2060 & $-\tfrac{1}{45}$ & $(\tfrac{1}{2},0,0)$ & \texttt{344477}$_{480}$ & 1751 & $\tfrac{791}{7200}$ & $(0,-1,0)$\\
\texttt{344666}$_{240}$ & 2748 & $-\tfrac{2}{75}$ & $(\tfrac{1}{2},0,0)$ & \texttt{344667}$_{240}$ & 2392 & $\tfrac{133}{1800}$ & $(0,-1,0)$\\
\texttt{344673}$_{480}$ & 2074 & $\tfrac{77}{3600}$ & 0 & \texttt{344676}$_{480}$ & 2566 & $\tfrac{173}{3600}$ & $(-\tfrac{1}{2},0,0)$\\
\texttt{344775}$_{240}$ & 1400 & $\tfrac{5}{72}$ & $(-\tfrac{1}{2},0,0)$ & \texttt{344777}$_{240}$ & 2018 & $-\tfrac{491}{3600}$ & $(0,1,\tfrac{1}{2})$\\
\texttt{346474}$_{480}$ & 1920 & $0$ & $(-\tfrac{1}{2},1,0)$ & \texttt{346475}$_{480}$ & 1766 & $\tfrac{103}{3600}$ & 0\\
\texttt{346477}$_{480}$ & 2155 & $\tfrac{11}{1440}$ & 0 & \texttt{346644}$_{240}$ & 2526 & $-\tfrac{79}{1200}$ & $(1,0,0)$\\
\texttt{346645}$_{240}$ & 2268 & $\tfrac{3}{200}$ & 0 & \texttt{346673}$_{240}$ & 2950 & $\tfrac{7}{720}$ & 0\\
\texttt{346674}$_{480}$ & 2774 & $\tfrac{157}{3600}$ & $(-\tfrac{1}{2},0,0)$ & \texttt{346744}$_{480}$ & 2116 & $\tfrac{229}{1800}$ & $(-\tfrac{1}{2},-1,0)$\\
\texttt{346774}$_{480}$ & 2395 & $-\tfrac{5}{288}$ & $(\tfrac{1}{2},0,0)$ & \texttt{346777}$_{480}$ & 2688 & $\tfrac{23}{200}$ & $(0,0,-\tfrac{1}{2})$\\
\texttt{347666}$_{480}$ & 3158 & $\tfrac{19}{3600}$ & 0 & \texttt{347667}$_{480}$ & 2887 & $\tfrac{67}{7200}$ & 0\\
\texttt{347745}$_{240}$ & 1564 & $\tfrac{91}{1800}$ & 0 & \texttt{347747}$_{240}$ & 2194 & $\tfrac{77}{3600}$ & 0\\
\texttt{347764}$_{480}$ & 2310 & $\tfrac{19}{240}$ & $(0,-1,0)$ & \texttt{347767}$_{480}$ & 2579 & $\tfrac{119}{7200}$ & 0\\
\texttt{366776}$_{240}$ & 3532 & $\tfrac{13}{1800}$ & 0 & \texttt{366777}$_{240}$ & 3370 & $\tfrac{79}{720}$ & $(0,0,-\tfrac{1}{2})$\\
\texttt{367766}$_{240}$ & 3632 & $\tfrac{1}{225}$ & 0 & \texttt{367767}$_{240}$ & 3416 & $\tfrac{7}{900}$ & 0\\
\texttt{444444}$_{40}$ & 900 & $\tfrac{1}{8}$ & 0 & \texttt{444456}$_{480}$ & 1260 & $\tfrac{1}{8}$ & $(-1,0,0)$\\
\texttt{444457}$_{480}$ & 1020 & $\tfrac{1}{8}$ & $(\tfrac{1}{2},-1,0)$ & \texttt{444467}$_{480}$ & 1740 & $\tfrac{1}{8}$ & $(-\tfrac{1}{2},-1,0)$\\
\texttt{444555}$_{80}$ & 360 & $\tfrac{1}{8}$ & $(\tfrac{3}{2},0,0)$ & \texttt{444557}$_{240}$ & 840 & $\tfrac{1}{8}$ & $(-1,1,0)$\\
\texttt{444577}$_{240}$ & 1320 & $-\tfrac{1}{8}$ & $(\tfrac{1}{2},2,0)$ & \texttt{444666}$_{80}$ & 2520 & $\tfrac{1}{8}$ & $(-\tfrac{3}{2},0,0)$\\
\texttt{444667}$_{240}$ & 2280 & $-\tfrac{1}{8}$ & $(1,1,0)$ & \texttt{444677}$_{240}$ & 2040 & $\tfrac{1}{8}$ & $(\tfrac{1}{2},-2,0)$\\
\texttt{445544}$_{120}$ & 576 & $\tfrac{9}{50}$ & 0 & \texttt{445546}$_{240}$ & 1146 & $-\tfrac{29}{1200}$ & $(1,0,0)$\\
\texttt{445577}$_{240}$ & 1240 & $\tfrac{29}{360}$ & $(-\tfrac{1}{2},0,0)$ & \texttt{445644}$_{240}$ & 1152 & $\tfrac{3}{50}$ & 0\\
\texttt{445647}$_{480}$ & 1420 & $-\tfrac{19}{360}$ & $(1,0,0)$ & \texttt{445666}$_{240}$ & 2292 & $-\tfrac{7}{300}$ & $(\tfrac{1}{2},0,0)$\\
\texttt{445667}$_{480}$ & 1994 & $\tfrac{187}{3600}$ & $(-\tfrac{1}{2},0,0)$ & \texttt{445746}$_{480}$ & 1436 & $-\tfrac{31}{1800}$ & $(-\tfrac{1}{2},1,0)$\\
\texttt{445747}$_{480}$ & 1249 & $\tfrac{1189}{7200}$ & $(-\tfrac{1}{2},-1,0)$ & \texttt{445766}$_{240}$ & 2168 & $\tfrac{137}{1800}$ & $(0,-1,0)$\\
\texttt{445777}$_{240}$ & 1582 & $\tfrac{311}{3600}$ & $(0,1,-\tfrac{1}{2})$ & \texttt{446555}$_{240}$ & 708 & $\tfrac{1}{150}$ & $(\tfrac{1}{2},0,0)$\\
\texttt{446557}$_{480}$ & 1346 & $\tfrac{223}{3600}$ & $(-\tfrac{1}{2},0,0)$ & \texttt{446644}$_{120}$ & 2304 & $\tfrac{1}{50}$ & 0\\
\texttt{446645}$_{240}$ & 2166 & $\tfrac{101}{1200}$ & $(-1,0,0)$ & \texttt{446677}$_{240}$ & 2840 & $-\tfrac{11}{360}$ & $(\tfrac{1}{2},0,0)$\\
\texttt{446745}$_{480}$ & 1756 & $-\tfrac{131}{1800}$ & $(\tfrac{1}{2},1,0)$ & \texttt{446747}$_{480}$ & 2281 & $-\tfrac{179}{7200}$ & $(-\tfrac{1}{2},1,0)$\\
\texttt{446755}$_{240}$ & 1528 & $\tfrac{187}{1800}$ & $(0,-1,0)$ & \texttt{446777}$_{240}$ & 2398 & $-\tfrac{121}{3600}$ & $(0,-1,\tfrac{1}{2})$\\
\texttt{447477}$_{160}$ & 1725 & $\tfrac{1}{32}$ & 0 & \texttt{447555}$_{240}$ & 568 & $\tfrac{307}{1800}$ & $(-\tfrac{1}{2},0,0)$\\
\texttt{447557}$_{480}$ & 1107 & $\tfrac{89}{2400}$ & $(\tfrac{1}{2},0,0)$ & \texttt{447577}$_{240}$ & 1438 & $-\tfrac{91}{3600}$ & $(\tfrac{1}{2},0,0)$\\
\texttt{447666}$_{240}$ & 2872 & $\tfrac{73}{1800}$ & $(-\tfrac{1}{2},0,0)$ & \texttt{447667}$_{480}$ & 2763 & $-\tfrac{79}{2400}$ & $(\tfrac{1}{2},0,0)$\\
\texttt{447677}$_{240}$ & 2498 & $\tfrac{199}{3600}$ & $(-\tfrac{1}{2},0,0)$ & \texttt{447744}$_{120}$ & 1600 & $\tfrac{1}{18}$ & 0\\
\texttt{447747}$_{240}$ & 1920 & $-\tfrac{1}{10}$ & $(0,2,0)$ & \texttt{447757}$_{480}$ & 1809 & $\tfrac{203}{2400}$ & $(0,-1,0)$\\
\texttt{447767}$_{480}$ & 2361 & $-\tfrac{133}{2400}$ & $(0,1,0)$ & \texttt{447777}$_{240}$ & 2160 & $\tfrac{1}{5}$ & $(0,-1,-\tfrac{1}{2})$\\
\texttt{455677}$_{480}$ & 1554 & $\tfrac{13}{400}$ & 0 & \texttt{456774}$_{480}$ & 1975 & $\tfrac{23}{1440}$ & 0\\
\texttt{457667}$_{480}$ & 2454 & $\tfrac{3}{400}$ & 0 & \texttt{457677}$_{480}$ & 2315 & $\tfrac{19}{1440}$ & 0\\
\texttt{457755}$_{240}$ & 1078 & $\tfrac{299}{3600}$ & 0 & \texttt{457775}$_{480}$ & 1509 & $\tfrac{83}{2400}$ & 0\\
\texttt{457777}$_{480}$ & 1937 & $-\tfrac{583}{7200}$ & $(0,0,\tfrac{1}{2})$ & \texttt{467766}$_{240}$ & 3382 & $\tfrac{11}{3600}$ & 0\\
\texttt{467776}$_{480}$ & 2979 & $\tfrac{13}{2400}$ & 0 & \texttt{467777}$_{480}$ & 2887 & $-\tfrac{713}{7200}$ & $(0,0,\tfrac{1}{2})$\\
\texttt{477777}$_{240}$ & 2550 & $\tfrac{1}{48}$ & 0 & \texttt{555555}$_{20}$ & 24 & $\tfrac{191}{300}$ & 0\\
\texttt{555557}$_{120}$ & 304 & $\tfrac{109}{450}$ & 0 & \texttt{555577}$_{240}$ & 668 & $\tfrac{137}{1800}$ & 0\\
\texttt{555777}$_{80}$ & 1146 & $-\tfrac{13}{400}$ & $(0,0,\tfrac{1}{2})$ & \texttt{557777}$_{240}$ & 1662 & $\tfrac{49}{400}$ & $(0,0,-\tfrac{1}{2})$\\
\texttt{577776}$_{120}$ & 2480 & $\tfrac{1}{90}$ & 0 & \texttt{577777}$_{120}$ & 2404 & $\tfrac{391}{1800}$ & $(0,0,-1)$\\
\texttt{666666}$_{20}$ & 4584 & $\tfrac{1}{300}$ & 0 & \texttt{666667}$_{120}$ & 4336 & $\tfrac{1}{450}$ & 0\\
\texttt{666677}$_{240}$ & 3988 & $\tfrac{7}{1800}$ & 0 & \texttt{666777}$_{80}$ & 3594 & $\tfrac{43}{400}$ & $(0,0,-\tfrac{1}{2})$\\
\texttt{667777}$_{240}$ & 3438 & $-\tfrac{39}{400}$ & $(0,0,\tfrac{1}{2})$ & \texttt{677777}$_{120}$ & 3164 & $\tfrac{401}{1800}$ & $(0,0,-1)$\\
\texttt{777777}$_{20}$ & 2880 & $-\tfrac{2}{5}$ & $(0,0,2)$ &  &  &  & \\
\bottomrule
\end{longtable}}

